\section{The formal content of a \safetyc{} comment}
\label{sec:theory}

Before we can evaluate designs, we must pin down the mathematical content of
the thing being proposed. This section fixes terminology used in the rest of
the study: what proposition an \code{unsafe} block's safety comment actually
denotes, what a proof of it would range over, and where Curry--Howard enters.

\subsection{The proposition behind an \code{unsafe} block}
\label{sec:theory-proposition}

Rust's soundness story is often summarized as: \emph{safe code cannot cause
undefined behavior}. The precise version, due to
RustBelt~\cite{rustbelt2018}, is a \emph{semantic typing} statement. Each
type $\tau$ is interpreted not as a syntactic classifier but as a predicate
$\sem{\tau}$ on values --- in RustBelt's case, a predicate in the Iris
separation logic~\cite{iris2018} over a low-level operational semantics of
MIR. A function is \emph{semantically well-typed} when it maps arguments
satisfying its parameter predicates to results satisfying its return
predicate, without ever executing an undefined operation, \emph{regardless
of what well-typed code surrounds it}. The fundamental theorem says
syntactically well-typed (safe) code is semantically well-typed; soundness
of a program containing \code{unsafe} then reduces to a sequence of manual
lemmas:

\begin{definition}[Safety obligation of an unsafe block]
\label{def:obligation}
Let $B$ be an \code{unsafe} block occurring in a function $f$ of module $M$,
and let $I_M$ be the module invariant --- the (possibly ghost-state-indexed)
predicate relating $M$'s private data structures to the abstraction $M$
exposes. The safety obligation of $B$ is the Hoare-style statement
\[
  \hoare{\,P_B \sep I_M\,}{\;B\;}{\,Q_B \sep I_M\,}
\]
where $P_B$ is what the surrounding safe code is known to have established
at entry to $B$ (from types, flow, and prior checks), and $Q_B$ is what the
code after $B$ relies on. ``Proving the SAFETY comment'' means exhibiting a
derivation of this triple against a formal operational model of the
language, in which every primitive step of $B$ is shown to satisfy its
side conditions (validity, alignment, aliasing, initialization, \dots).
\end{definition}

Three consequences of this framing structure the rest of the study.

\paragraph{(1) The obligation is relative to \emph{two} semantics.}
The side conditions in Definition~\ref{def:obligation} come from Rust's
dynamic semantics --- pointer validity and provenance as fixed by the
operational aliasing model (Stacked Borrows~\cite{stackedborrows2019}, now
superseded in practice by Tree Borrows~\cite{treeborrows2025}) --- but the
\emph{preconditions} $P_B$ in \dotnetrs{} are frequently facts about the
\emph{virtual machine being implemented}, not about Rust: ``all mutator
threads are stopped,'' ``this object's arena is registered and its
generation matches,'' ``the layout descriptor says offset $k$ holds an
object reference.'' Those facts are only expressible against a model of the
VES itself --- an abstract machine for a fragment of
ECMA-335~\cite{ecma335}. A verifier for \dotnetrs{} safety comments is
therefore necessarily a \emph{two-level} system:
\[
  \underbrace{\text{ECMA-335 abstract machine model}}_{\text{gives meaning to } P_B, Q_B, I_M}
  \quad\Longrightarrow\quad
  \underbrace{\text{Rust/MIR memory model}}_{\text{gives meaning to the steps of } B}
\]
with a simulation relation connecting VM-level states to the Rust data
structures that represent them. This is the same architecture as a compiler
correctness proof (CompCert~\cite{compcert2009}, CakeML~\cite{cakeml2014}),
run in reverse: instead of proving a compiler preserves source semantics, we
prove a hand-written interpreter's internal transitions refine the
specification's transitions, and that its \code{unsafe} shortcuts are
justified by invariants the specification-level model maintains.

\paragraph{(2) Module invariants, not block-local facts, carry the weight.}
Almost none of the interesting \safetyc{} comments in \dotnetrs{} are
self-contained. A comment like ``STW has stopped every mutator, so the
owning arena cannot be freed during this read''
(\cref{sec:obligations}) appeals to a \emph{protocol}: a
state machine over coordinator phases, arena registration, lease counts,
and thread states, whose transitions are scattered across several crates.
In separation-logic terms this is ghost state governed by an invariant
$I_{\mathrm{STW}}$; individual blocks open the invariant, use a token
certifying the current phase, and close it. Any DSL that can only state
block-local, first-order facts (``$p$ is aligned,'' ``$n \le \mathrm{len}$'')
will be able to discharge the easy third of the corpus and will hit a wall
exactly where the risk is concentrated. This is why the study devotes
\cref{sec:design} to the question of \emph{where the invariant layer
lives}, and it is the single strongest argument for reusing an existing
Iris-class logic rather than inventing a bespoke solver.

\paragraph{(3) The proof is about a model of the code, not the code.}
Unless the checker consumes the actual MIR of the crate (as
RefinedRust~\cite{refinedrust2024} and Aeneas~\cite{aeneas2022} do), there
is a formalization gap: someone (or some tool) must translate each
\code{unsafe} block and the definitions it touches into the proof
language, and the fidelity of that translation is part of the trusted
computing base. \Cref{sec:risks} quantifies this: for a codebase undergoing
LLM-assisted refactors at \dotnetrs{}'s pace, an untethered model drifts
from the code within weeks. Mechanized tethering --- extraction from MIR, or
at minimum syntactic anchoring with CI-enforced drift checks in the style of
the project's existing \code{check\_doc\_drift.sh} --- is not optional
polish; it is what distinguishes verification from documentation.

\subsection{The role of the Curry--Howard correspondence}
\label{sec:theory-ch}

The Curry--Howard correspondence identifies propositions with types and
proofs with programs: a term $t : \Phi$ \emph{is} a checkable certificate of
$\Phi$. For this project the correspondence matters in a specific,
practical way: it determines the \emph{artifact format}. A
\safetyc{}-comment DSL under Curry--Howard is a language of \emph{proof
terms} whose types are safety propositions; the checker is a type checker;
and soundness of the whole edifice reduces to (a) soundness of the ambient
type theory and (b) adequacy of the axioms/model against which the
propositions are stated. Concretely, three design families instantiate
this:

\begin{description}
\item[Deep proof terms, custom checker.] The DSL is itself a small
  dependently-typed calculus ($\Pi$, $\Sigma$, indexed inductive families
  for the invariant taxonomy of \cref{sec:obligations}), with a
  hand-written kernel. Curry--Howard is used directly. The kernel joins
  the TCB; decidability and totality of type checking must be engineered
  (bidirectional checking, no unrestricted conversion); and every
  automation convenience (arithmetic, congruence) must be built, or
  the proofs become impractically manual. Bespoke kernels developed
  alongside large host projects have a poor historical record;
  \cref{sec:design-custom} argues this option is dominated.
\item[Shallow embedding into an existing proof assistant.] Propositions
  and proofs are ordinary Lean~4~\cite{lean4} or Rocq~\cite{rocq}
  terms; the annotation beside the \code{unsafe} site --- whether an
  attribute, a comment, or an embedded proof script --- determines the
  statement and sketches the proof, while the checked artifact lives
  in the host. Curry--Howard is
  inherited from the host's kernel, which is small, battle-tested, and
  independently checkable. The VM model, the Rust memory-model fragment,
  and the simulation relation are libraries in the host. This is the
  architecture of RustBelt (Iris/Rocq), RefinedRust (Iris/Rocq, with a
  MIR front end), and Aeneas (extraction to Lean).
\item[SMT-mediated verification conditions.] Verus~\cite{verus2024} and
  Creusot~\cite{creusot2022} compile annotations to SMT queries. This
  buys enormous automation for first-order obligations (bounds,
  alignment, arithmetic) at the cost of the proof-term ideal: an SMT
  ``unsat'' is not a checkable certificate in the Curry--Howard sense,
  and quantifier-heavy protocol invariants strain the paradigm. The
  pragmatic middle ground --- proof terms for protocol lemmas, SMT for
  leaf arithmetic --- is exactly the division of labor RefinedRust and
  Verus each approach from opposite ends.
\end{description}

\subsection{The specification side: an ECMA-335 abstract machine}
\label{sec:theory-vesmodel}

The second layer --- the one that distinguishes this proposal from
general unsafe-Rust verification --- is a formal model of the VES fragment that \dotnetrs{}
implements. The standard itself is prose, but it is unusually
formalization-friendly prose: it specifies an explicit evaluation-stack
machine, a typed memory store with a precise alignment and atomicity
contract (\ecma{I.12.6.2}, \ecma{I.12.6.6}), an acquire/release volatile
semantics (\ecma{I.12.6.7}), a pointer taxonomy separating managed from
unmanaged pointers with GC-reporting obligations (\ecma{II.14.4},
\ecma{III.1.1.5}), layout rules including the prohibition on object
references overlapping other fields (\ecma{II.10.7}), and finalization
lifecycle rules (\ecma{I.8.9.6.7}). Prior mechanizations of comparable
scope exist and calibrate the effort: Fruja's ASM-based CLR model covering
the type system and bytecode verifier~\cite{fruja2006}, Gordon and Syme's
type soundness for a CIL fragment~\cite{gordonsyme2001}, Kennedy and Syme's
formalization of CLR generics~\cite{kennedysyme2001}, and --- closest in
spirit to what \dotnetrs{} needs --- the Jinja/JinjaThreads line in
Isabelle/HOL, which grew from a sequential Java-like machine to a verified
multithreaded JVM with a verified compiler over roughly a
decade~\cite{jinja2006,jinjathreads2012}.

The model does \emph{not} need to cover the whole standard to be useful.
The invariant families of \cref{sec:obligations} touch a compact core:
object and array layout, the managed-pointer discipline, the memory/atomicity
model, GC reachability and finalization, and the STW protocol (which is an
implementation construct \emph{justified by} the standard's GC latitude, not
mandated by it). A model of that core, indexed by the layout function that
\dotnetrs{}'s \code{LayoutManager} actually computes, is the central
artifact; everything else in the roadmap of \cref{sec:roadmap} is staged
against it.
